% Part: many-valued-logic
% Chapter: syntax-and-semantics
% Section: semantic-notions

\documentclass[../../../include/open-logic-section]{subfiles}

\begin{document}

\olfileid{mvl}{syn}{sem}

\olsection{અર્થવિચારી ખ્યાલો}

ધારો કે બહુમૂલ્ય તર્કશાસ્ત્ર~$\Log L$ શ્રેણિક વડે આપેલું છે. ત્યારે
આપણે~$\Log L$ માટેના સામાન્ય અર્થવિચારી ખ્યાલો વ્યાખ્યાયિત કરી શકીએ.

\begin{defn}
\begin{enumerate}
\item કોઈ~$\pAssign{v}$ માટે $\pSat{v}{!A}$ હોય, તો
  !!^a{formula}~$!A$ \emph{સંતોષ્ય} છે; કોઈ પણ $\pAssign{v}$ માટે
  $\pSat{v}{!A}$ ન હોય, તો તે \emph{અસંતોષ્ય} છે;
\item બધાં !!{valuation}s~$v$ માટે $\pSat{v}{!A}$ હોય, તો
  !!^a{formula}~$!A$ \emph{પુનરુક્તિ} છે;
\item જો $\Gamma$ એ !!{formula}sનો ગણ હોય, તો $\Gamma \Entails !A$
  (``$\Gamma$, $!A$ને ફલિત કરે છે'') ત્યારે અને માત્ર ત્યારે, જ્યારે
  $\pSat{v}{\Gamma}$ હોય એવા દરેક !!{valuation}~$\pAssign{v}$ માટે
  $\pSat{v}{!A}$ હોય.
\item જો $\Gamma$ એ !!{formula}sનો ગણ હોય, તો એવા કોઈ
  !!a{valuation}~$\pAssign{v}$ માટે $\pSat{v}{\Gamma}$ હોય ત્યારે
  $\Gamma$ \emph{સંતોષ્ય} છે; અન્યથા $\Gamma$ \emph{અસંતોષ્ય} છે.
\end{enumerate}
\end{defn}

આ ખ્યાલો માટે શાસ્ત્રીય તર્કશાસ્ત્રના કિસ્સા જેવી કેટલીક હકીકતો મળે છે:

\begin{prop}
\ollabel{prop:semanticalfacts}
\begin{enumerate}
\item $!A$ પુનરુક્તિ હોય ત્યારે અને માત્ર ત્યારે
  $\emptyset \Entails !A$;
\item જો $\Gamma$ સંતોષ્ય હોય, તો $\Gamma$નો દરેક સાન્ત ઉપગણ પણ
  સંતોષ્ય છે;
\item\ollabel{def:monotonicity}%
એકદિશવર્ધિતા: જો $\Gamma \subseteq \Delta$
  અને $\Gamma \Entails !A$, તો $\Delta \Entails !A$ પણ;
\item\ollabel{def:Cut}%
પરંપરિતતા: જો $\Gamma \Entails !A$ અને
  $\Delta \cup \{ !A\} \Entails !B$, તો $\Gamma \cup \Delta \Entails
  !B$;
\end{enumerate}
\end{prop}

\begin{proof}
અભ્યાસપ્રશ્ન.
\end{proof}

\begin{prob}
\olref[mvl][syn][sem]{prop:semanticalfacts} સાબિત કરો.
\end{prob}

શાસ્ત્રીય તર્કશાસ્ત્રમાં ફલિતતા અને પ્રેરણ સંયોજક વચ્ચે સંબંધ સ્થાપી શકાય છે.
દાખલા તરીકે, \emph{મોદુસ પોનેન્સ}નું પ્રામાણ્ય મળે છે: જો $\Gamma
\Entails !A$ અને $\Gamma \Entails !A \lif !B$, તો $\Gamma \Entails
!B$. શાસ્ત્રીય તર્કશાસ્ત્રમાં $\Entails$ અને $\lif$ વચ્ચેનો બીજો
મહત્ત્વનો સંબંધ અર્થાનુસારી નિગમન પ્રમેય છે: $\Gamma \Entails !A
\lif !B$ ત્યારે અને માત્ર ત્યારે, જ્યારે $\Gamma \cup \{!A\} \Entails
!B$. આ પરિણામો બહુમૂલ્ય તર્કશાસ્ત્રોમાં \emph{હંમેશાં લાગુ પડતાં નથી}.
તેઓ લાગુ પડે છે કે કેમ તે સત્યવિધેય~$\tf{\lif}$ પર આધાર રાખે છે.

\end{document}
